% DO NOT EDIT - automatically generated from metadata.yaml

\def \codeURL{https://github.com/DovydasVe/topological-loss-rescience/}
\def \codeDOI{}
\def \dataURL{}
\def \dataDOI{}
\def \editorNAME{}
\def \editorORCID{}
\def \reviewerINAME{}
\def \reviewerIORCID{}
\def \reviewerIINAME{}
\def \reviewerIIORCID{}
\def \dateRECEIVED{}
\def \dateACCEPTED{}
\def \datePUBLISHED{}
\def \articleTITLE{[Re] Topological Loss Function for Image Segmentation}
\def \articleTYPE{Replication}
\def \articleDOMAIN{Segmentation}
\def \articleBIBLIOGRAPHY{bibliography.bib}
\def \articleYEAR{2026}
\def \reviewURL{}
\def \articleABSTRACT{We present an independent replication of Clough et al. (2022), ``A Topological Loss Function for Deep-Learning Based Image Segmentation Using Persistent Homology.'' We implement the topological post-processing framework from scratch and assess its performance across Experiment 1 (MNIST image denoising) and Experiment 3 (ACDC cardiac MRI myocardium segmentation). Our results confirm statistically significant classification accuracy improvements across all noise levels on MNIST and modest volumetric 3D Dice improvements (+0.41 percentage points on end-diastole and +0.30 percentage points on end-systole) on ACDC.}
\def \replicationCITE{J. R. Clough, N. Byrne, I. Oksuz, V. A. Zimmer, J. A. Schnabel and A. P. King, ``A Topological Loss Function for Deep-Learning Based Image Segmentation Using Persistent Homology,'' in IEEE Transactions on Pattern Analysis and Machine Intelligence, vol. 44, no. 12, pp. 8766-8778, 1 Dec. 2022.}
\def \replicationBIB{Clough:2022}
\def \replicationURL{https://ieeexplore.ieee.org/document/9186664/}
\def \replicationDOI{10.1109/TPAMI.2020.3013679}
\def \contactNAME{Dovydas Venclovas}
\def \contactEMAIL{dhvenclovas@gmail.com}
\def \articleKEYWORDS{rescience c, python, topological data analysis, topological loss function, neural networks, deep learning}
\def \journalNAME{ReScience C}
\def \journalVOLUME{}
\def \journalISSUE{}
\def \articleNUMBER{}
\def \articleDOI{}
\def \authorsFULL{Dovydas Venclovas}
\def \authorsABBRV{D. Venclovas}
\def \authorsSHORT{Venclovas}
\title{\articleTITLE}
\date{}
\author[1\orcid{0009-0006-1523-3853}]{Dovydas Venclovas}
\affil[1]{University of Amsterdam, Amsterdam, The Netherlands}